\chapter{المتجهات والمستقيمات في المستوي}\label{ch:g11:vect}

تشفّر المتجهات الانتقالات؛ والمستقيمات مساراتها. والأداة الجديدة
الأساسية في هذا الفصل هي \emph{\omterm{def:g11:vect:det}{المحدد}}، وهو عدد مفرد
يقرر هل \omterm{def:g11:vect:collinear}{متجهتان مستقيميتان} --- ومن هناك ينتج
معادلات المستقيمات ويحسم مسائل الاستقامة والتوازي بالحساب
الخالص. وتمتد الأفكار نفسها إلى الفضاء في \cref{ch:g12:space}.

\section{المتجهات: تذكير}

\begin{definition}[المتجهة والإحداثيات]\label{def:g11:vect:vector}
تمثّل \emph{المتجهة}\index{متجهة} $\vec u$ انتقالًا: فكل
الأسهم ذات الاتجاه والجهة والطول نفسها تمثّل المتجهة
نفسها. وفي \omterm{def:g10:coordgeom:system}{نظام إحداثيات}، تسجّل $\vec u\,(x, y)$ الانتقال
($x$ أفقيًا و $y$ شاقوليًا)؛ ومن أجل النقطتين $A(x_A, y_A)$ و
$B(x_B, y_B)$،
\[
\vect{AB}\,\bigl(x_B - x_A,\ y_B - y_A\bigr).
\]
وتُجمع المتجهات إحداثيًا إحداثيًا، وللمتجهة $\lambda\vec u$ \omterm{def:g10:coordgeom:system}{الإحداثيات}
$(\lambda x, \lambda y)$.
\end{definition}

\begin{proposition}[المنتصف]\label{prop:g11:vect:midpoint}
\omterm{prop:g10:coordgeom:midpoint}{لمنتصف} القطعة $[AB]$، وليكن $M$، \omterm{def:g10:coordgeom:system}{الإحداثيات}
$\left(\frac{x_A + x_B}{2},\ \frac{y_A + y_B}{2}\right)$.
\end{proposition}

\begin{proof}
تتميز $M$ بالعلاقة $\vect{AM} = \frac12\vect{AB}$: وإحداثيًا إحداثيًا،
$x_M - x_A = \frac12(x_B - x_A)$، إذن $x_M = \frac{x_A + x_B}{2}$، و
بالمثل من أجل $y_M$.
\end{proof}

\section{الاستقامية والمحدد}

\begin{definition}[المتجهتان المستقيميتان]\label{def:g11:vect:collinear}
تكون متجهتان \emph{مستقيميتين}\index{متجهتان مستقيميتان} إذا كانت إحداهما
مضاعفًا سلّميًا للأخرى: $\vec v = \lambda \vec u$ من أجل
$\lambda \in \R$ ما (أو $\vec u = \vec 0$). وهندسيًا: لهما الاتجاه
نفسه، أو إحداهما معدومة.
\end{definition}

\begin{definition}[المحدد]\label{def:g11:vect:det}
\emph{محدد}\index{محدد} المتجهتين $\vec u\,(x, y)$ و
$\vec v\,(x', y')$ هو
\[
\det(\vec u, \vec v) =
\begin{vmatrix} x & x' \\ y & y' \end{vmatrix}
= x y' - x' y .
\]
\end{definition}

\begin{theorem}[محك الاستقامية]\label{thm:g11:vect:collinearity}
تكون المتجهتان $\vec u$ و $\vec v$ \omterm{def:g11:vect:collinear}{مستقيميتين} إذا وفقط إذا كان
$\det(\vec u, \vec v) = 0$.
\end{theorem}

\begin{proof}
($\Rightarrow$) إذا كانت $\vec v = \lambda\vec u$، فإن $x' = \lambda x$ و
$y' = \lambda y$، إذن
$\det(\vec u, \vec v) = x(\lambda y) - (\lambda x)y = 0$؛ وإذا كانت
$\vec u = \vec 0$، \omterm{def:g11:vect:det}{فالمحدد} معدوم بداهةً.

($\Leftarrow$) لنفترض أن $xy' - x'y = 0$ و $\vec u \neq \vec 0$، وليكن مثلًا
$x \neq 0$ (والحالة $y \neq 0$ متناظرة). ضع
$\lambda = \frac{x'}{x}$. عندئذٍ $x' = \lambda x$ بالإنشاء، وتعطي
العلاقة $xy' = x'y = \lambda x y$، بعد القسمة على $x \neq 0$،
أن $y' = \lambda y$. ومنه $\vec v = \lambda \vec u$.
\end{proof}

\begin{example}\label{ex:g11:vect:collinear}
$\vec u\,(3, -2)$ و $\vec v\,(-6, 4)$:
$\det = 3 \times 4 - (-6) \times (-2) = 12 - 12 = 0$، إذن مستقيميتان (وفعلًا
$\vec v = -2\vec u$). و $\vec u\,(3, -2)$ و $\vec w\,(2, 1)$:
$\det = 3 + 4 = 7 \neq 0$، فليستا \omterm{def:g11:vect:collinear}{مستقيميتين}.
\end{example}

\begin{method}[استقامة ثلاث نقط]\label{met:g11:vect:alignment}
تكون $A$ و $B$ و $C$ على استقامة واحدة إذا وفقط إذا كانت $\vect{AB}$ و $\vect{AC}$
\omterm{def:g11:vect:collinear}{مستقيميتين}: فاحسب المتجهتين وتحقق من
$\det\bigl(\vect{AB}, \vect{AC}\bigr) = 0$.
\end{method}

\section{المعادلات الديكارتية للمستقيمات}

\begin{definition}[متجهة التوجيه]\label{def:g11:vect:direction}
\emph{متجهة التوجيه}\index{متجهة توجيه} لمستقيم $d$ هي أي
\omterm{def:g11:vect:vector}{متجهة} غير معدومة $\vec u$ تكون $\vect{AB}$ مستقيمية مع $\vec u$
من أجل كل نقطتين $A, B$ من $d$.
\end{definition}

\begin{theorem}[المعادلة الديكارتية]\label{thm:g11:vect:cartesian}
للمستقيم ذي \omterm{def:g11:vect:direction}{متجهة التوجيه} $\vec u\,(-b, a)$ \omterm{def:g10:algebra:equation}{معادلة} على
الصورة
\[
ax + by + c = 0
\qquad (a, b) \neq (0, 0).
\]
وبالعكس، تصف كل \omterm{def:g10:algebra:equation}{معادلة} كهذه مستقيمًا \omterm{def:g11:vect:vector}{متجهة} توجيهه
$\vec u\,(-b, a)$.
\end{theorem}

\begin{proof}
ليمر $d$ بالنقطة $A(x_A, y_A)$ وليكن توجيهه $\vec u\,(-b, a)$. تقع
النقطة $M(x, y)$ على $d$ بالضبط عندما تكون $\vect{AM}$ مستقيمية مع
$\vec u$، أي (\cref{thm:g11:vect:collinearity})
\[
0 = \det\bigl(\vect{AM}, \vec u\bigr)
= (x - x_A)\,a - (-b)(y - y_A)
= ax + by - (a x_A + b y_A),
\]
وهي \omterm{def:g10:algebra:equation}{معادلة} على الصورة المعلنة مع $c = -(a x_A + b y_A)$.
وبالعكس، فإن حلول $ax + by + c = 0$ مع، مثلًا، $b \neq 0$ هي
النقط $\left(x, -\frac{a x + c}{b}\right)$: فطرح اثنتين منها
يبيّن أن كل وتر مستقيمي مع $(-b, a)$، ويوجد حل دائمًا
--- فهو المستقيم المار به والموجَّه \omterm{def:g11:vect:vector}{بالمتجهة} $(-b, a)$.
\end{proof}

\begin{method}[المستقيم المار بنقطتين]\label{met:g11:vect:twopoints}
لإيجاد \omterm{def:g10:algebra:equation}{معادلة} المستقيم $(AB)$: احسب \omterm{def:g11:vect:direction}{متجهة التوجيه}
$\vect{AB}\,(x_B - x_A, y_B - y_A)$؛ وطابقها مع $(-b, a)$، أي خذ
$a = y_B - y_A$ و $b = -(x_B - x_A)$؛ ثم حدّد $c$ بتعويض
\omterm{def:g10:coordgeom:system}{إحداثيات} $A$.
\end{method}

\begin{example}\label{ex:g11:vect:twopoints}
المستقيم المار بالنقطتين $A(1, 2)$ و $B(4, 3)$: $\vect{AB}\,(3, 1)$، إذن
$a = 1$ و $b = -3$، \omterm{def:g10:algebra:equation}{والمعادلة} $x - 3y + c = 0$. وبتعويض $A$:
$1 - 6 + c = 0$، إذن $c = 5$. \omterm{def:g10:algebra:equation}{والمعادلة}: $x - 3y + 5 = 0$. وتحقق بالنقطة $B$:
$4 - 9 + 5 = 0$.
\end{example}

\begin{remark}
عندما $b \neq 0$، يمكن حل \omterm{def:g10:algebra:equation}{المعادلة} بالنسبة إلى $y$: فتصير $y = mx + p$ حيث
\omterm{def:g10:reffunc:affine}{الميل} $m = -\frac ab$. والصورة الديكارتية $ax + by + c = 0$ أعم:
فهي تغطي أيضًا المستقيمات الشاقولية ($b = 0$) التي لا \omterm{def:g10:reffunc:affine}{ميل} لها.
\end{remark}

\begin{omfigure}
\begin{tikzpicture}
\begin{axis}[omaxis,
  width=8.6cm, height=5.8cm,
  xmin=-1.6, xmax=5.6, ymin=-0.9, ymax=4.4,
  xtick={1,4}, ytick={2,3}]
  \addplot[omDef, thick, domain=-1.4:5.4] {(x+5)/3};
  \fill (axis cs:1,2) circle (1.5pt)
    node[above left, font=\small] {$A$};
  \fill (axis cs:4,3) circle (1.5pt)
    node[above left, font=\small] {$B$};
  \draw[-Stealth, omThm, thick] (axis cs:1,2) -- (axis cs:2.5,2.5)
    node[midway, below right, font=\small] {$\vec u$};
\end{axis}
\end{tikzpicture}

{\small المستقيم $x - 3y + 5 = 0$ في \cref{ex:g11:vect:twopoints}، \omterm{def:g11:vect:vector}{ومتجهة}
توجيهه $\vec u = \vect{AB}$ (بالأحمر، وهي هنا مضروبة في $\frac12$)،
والنقطتان المستعملتان لبناء \omterm{def:g10:algebra:equation}{المعادلة}.}
\end{omfigure}

\section{التمثيل الوسيطي}

\begin{definition}[التمثيل الوسيطي]\label{def:g11:vect:parametric}
المستقيم المار بالنقطة $A(x_A, y_A)$ وتوجيهه $\vec u\,(\alpha, \beta)$ هو
مجموعة النقط
\[
M\bigl(x_A + t\alpha,\ y_A + t\beta\bigr), \qquad t \in \R .
\]
والعدد $t$ هو \emph{\omterm{def:g10:stats:median}{الوسيط}}: فكل قيمة له تنتج نقطة
واحدة من المستقيم، كأننا نسير عليه بسرعة ثابتة $\vec u$.
\end{definition}

\begin{example}
المستقيم في \cref{ex:g11:vect:twopoints} هو أيضًا
$\{(1 + 3t,\ 2 + t) : t \in \R\}$: إذ $t = 0$ يعطي $A$، و $t = 1$ يعطي $B$.
وبحذف $t$ (إذ $t = y - 2$، فيكون $x = 1 + 3(y-2)$) نستعيد
$x - 3y + 5 = 0$.
\end{example}

\section{الأوضاع النسبية لمستقيمين}

\begin{proposition}[متوازيان أم متقاطعان]\label{prop:g11:vect:positions}
يكون مستقيمان متجهتا توجيههما $\vec u$ و $\vec v$ متوازيين إذا وفقط
إذا كان $\det(\vec u, \vec v) = 0$. وبصيغة المعادلات: يكون $d : ax + by + c =
0$ و $d' : a'x + b'y + c' = 0$ متوازيين إذا وفقط إذا كان
$ab' - a'b = 0$. أما المستقيمان غير المتوازيين فيلتقيان في نقطة واحدة بالضبط.
\end{proposition}

\begin{proof}
التوازي يعني استقامية التوجيهين، وهذا هو
\cref{thm:g11:vect:collinearity}؛ ومع $\vec u\,(-b, a)$ و
$\vec v\,(-b', a')$، يكون \omterm{def:g11:vect:det}{المحدد} $(-b)a' - (-b')a = ab' - a'b$.
وإذا لم يكن المستقيمان متوازيين، فإن حل المعادلتين في آنٍ واحد
(بالتعويض أو بالتركيب) يقود إلى حل وحيد: فالجملة
تتصرف كجملة $2 \times 2$ ذات \omterm{def:g11:vect:det}{محدد} غير معدوم.
\end{proof}

\begin{method}[تقاطع مستقيمين]\label{met:g11:vect:intersection}
حل جملة المعادلتين: اعزل مجهولًا في إحدى المعادلتين
وعوّضه في الأخرى، أو ركّب المعادلتين لحذف أحد
المجهولين. وتحقق دائمًا من النقطة الناتجة في \emph{كلتا} المعادلتين.
\end{method}

\begin{example}\label{ex:g11:vect:intersection}
\omterm{def:g11:deriv:rate}{قاطع} $d: x - 3y + 5 = 0$ و $d': 2x + y - 4 = 0$. من $d'$:
$y = 4 - 2x$. وبالتعويض في $d$:
\[
x - 3(4 - 2x) + 5 = 0
\iff x - 12 + 6x + 5 = 0
\iff 7x = 7
\iff x = 1,
\]
ثم $y = 2$. إذن يلتقي المستقيمان في $(1, 2)$. وتحقق في $d$: $1 - 6 + 5 = 0$.
\end{example}

\section{تمارين}

\begin{exercise}[$\star$]\label{exo:g11:vect:1}
إذا أُعطيت $A(2, -1)$ و $B(5, 3)$ و $C(-1, 1)$، احسب \omterm{def:g10:coordgeom:system}{إحداثيات}
$\vect{AB}$ و $\vect{AC}$ و $\vect{AB} + \vect{AC}$ و $2\vect{AB} -
3\vect{AC}$، \omterm{prop:g10:coordgeom:midpoint}{ومنتصف} $[BC]$.
\end{exercise}

\begin{exercise}[$\star$]\label{exo:g11:vect:2}
هل \omterm{def:g11:vect:collinear}{المتجهتان مستقيميتان}؟
\[
\vec u\,(4, -6) \text{ و } \vec v\,(-2, 3); \qquad
\vec u\,(2, 5) \text{ و } \vec v\,(3, 7); \qquad
\vec u\,(0, 3) \text{ و } \vec v\,(0, -8).
\]
\end{exercise}

\begin{exercise}[$\star$]\label{exo:g11:vect:3}
أوجد \omterm{def:g10:algebra:equation}{معادلة} ديكارتية للمستقيم المار بالنقطة $A(2, 1)$ وتوجيهه
$\vec u\,(3, -1)$، ثم للمستقيم المار بالنقطتين $B(0, 4)$ و
$C(2, 0)$.
\end{exercise}

\begin{exercise}[$\star$]\label{exo:g11:vect:4}
من أجل المستقيم $d: 3x - 2y + 6 = 0$، أعطِ \omterm{def:g11:vect:direction}{متجهة توجيه}، و
نقطتي \omterm{def:g10:numbers:interunion}{التقاطع} مع محورَي \omterm{def:g10:coordgeom:system}{الإحداثيات}، وقرّر هل تنتمي النقطتان
$P(2, 6)$ و $Q(1, 4)$ إلى $d$.
\end{exercise}

\begin{exercise}[$\star\star$]\label{exo:g11:vect:5}
هل النقط $A(1, 2)$ و $B(4, 8)$ و $C(-2, -4)$ على استقامة واحدة؟ والسؤال نفسه
من أجل $D(0, 1)$ و $E(2, 4)$ و $F(5, 9)$.
\end{exercise}

\begin{exercise}[$\star\star$]\label{exo:g11:vect:6}
حدّد نقطة \omterm{def:g10:numbers:interunion}{التقاطع} إن وُجدت:
\[
d: 2x + y - 7 = 0 \ \text{ و }\ d': x - y + 1 = 0;
\qquad
e: x - 2y + 3 = 0 \ \text{ و }\ e': -2x + 4y + 1 = 0 .
\]
\end{exercise}

\begin{exercise}[$\star\star$]\label{exo:g11:vect:7}
أعطِ تمثيلًا وسيطيًا للمستقيم $d: 2x - y + 3 = 0$، \omterm{def:g10:algebra:equation}{ومعادلة}
ديكارتية للمستقيم
$\bigl\{(1 + 2t,\ -3 + t) : t \in \R\bigr\}$.
\end{exercise}

\begin{exercise}[$\star\star$]\label{exo:g11:vect:8}
أوجد \omterm{def:g10:algebra:equation}{معادلة} المستقيم المار بالنقطة $P(1, -2)$ الموازي للمستقيم
$d: 4x - 3y + 2 = 0$، وقيمة $m$ التي يكون من أجلها المستقيم
$mx + 2y - 5 = 0$ موازيًا للمستقيم $d$.
\end{exercise}

\begin{exercise}[$\star\star$]\label{exo:g11:vect:9}
لتكن $A(-1, 0)$ و $B(3, 2)$ و $C(1, -4)$. أوجد \omterm{def:g10:algebra:equation}{معادلة} ديكارتية
\emph{\omterm{def:g10:stats:mean}{لمتوسط}} المثلث $ABC$ الصادر من $C$ (وهو المستقيم الواصل بين $C$
\omterm{prop:g10:coordgeom:midpoint}{ومنتصف} $[AB]$).
\end{exercise}

\begin{exercise}[$\star\star$]\label{exo:g11:vect:10}
من أجل أي قيمة أو قيم \omterm{def:g10:stats:median}{للوسيط} $m$ تكون المتجهتان
$\vec u\,(m, 2)$ و $\vec v\,(3, m - 1)$ \omterm{def:g11:vect:collinear}{مستقيميتين}؟
\end{exercise}

\begin{exercise}[$\star\star\star$]\label{exo:g11:vect:11}
ليكن $ABCD$ متوازي أضلاع، ولتكن $I$ \omterm{prop:g10:coordgeom:midpoint}{منتصف} $[AB]$، ولتكن $J$
النقطة المعرَّفة بالعلاقة $\vect{DJ} = \frac23 \vect{DI}$. وبالعمل في
\omterm{def:g10:coordgeom:system}{نظام الإحداثيات} حيث $A(0,0)$ و $B(1,0)$ و $D(0,1)$ (فيكون $C(1,1)$):
\begin{enumerate}
  \item احسب \omterm{def:g10:coordgeom:system}{إحداثيات} $I$ و $J$.
  \item بيّن أن $A$ و $J$ و $C$ على استقامة واحدة. (وواقعة لطيفة: يقطع المستقيم
        $(DI)$ القطر $(AC)$ على ثلث ارتفاعه.)
\end{enumerate}
\end{exercise}

\section{مسألة: مسارات التصادم والإحداثيات الذكية}

\begin{problem}[{مسألة نهاية الأسبوع --- تقاطع المسارات ليس
تصادمًا: أقرب اقتراب في البحر، ومبرهنات تُبرهن باختيار
المعلم المناسب}]
\label{pb:g11:vect:1}
يتقاطع مساران مستقيمان لسفينتين --- فهل يجب أن تتصادم السفينتان؟ لا،
إلا إذا بلغتا نقطة \omterm{def:g10:numbers:interunion}{التقاطع} \emph{في اللحظة نفسها}: فالمسارات
مجموعات نقط، والحركات وسيطية، وقد أغرق الخلط بينهما
سفنًا حقيقية. تدير هذه المسألة خدمة صغيرة لحركة السفن
على مستقيمات وسيطية (\cref{def:g11:vect:parametric})، ثم تنتقل
إلى الهندسة الخالصة مع القوة الخارقة الأخرى \omterm{def:g10:coordgeom:system}{للإحداثيات}:
\emph{اختيارها} بذكاء، كما في \cref{exo:g11:vect:11}، حتى
تحسب المبرهنات الكلاسيكية نفسها بنفسها.

\medskip
\noindent\textbf{الجزء الأول --- المستقيم بثلاث طرائق.}
\begin{enumerate}
  \item أعطِ تمثيلًا وسيطيًا للمستقيم المار بالنقطة
        $A(1, 2)$ وتوجيهه $\vec u(3, -1)$، ثم
        احذف \omterm{def:g10:stats:median}{الوسيط} للحصول على \omterm{def:g10:algebra:equation}{معادلة} ديكارتية.
  \item من أجل المستقيم $2x - 5y + 3 = 0$: اقرأ \omterm{def:g11:vect:direction}{متجهة توجيه}
        (\cref{thm:g11:vect:cartesian})، وأوجد نقطة واحدة،
        واكتب تمثيلًا وسيطيًا.
  \item احسب نقطة \omterm{def:g10:numbers:interunion}{تقاطع} مستقيمَي
        السؤالين 1 و 2 (\cref{met:g11:vect:intersection}).
  \item أكّد \omterm{def:g11:vect:det}{بالمحدد}
        (\cref{thm:g11:vect:collinearity}) أن هذين المستقيمين
        كان لا بد أن يتقاطعا؛ ثم بيّن أن $2x - 5y + 3 = 0$ و
        $-4x + 10y + 1 = 0$ متوازيان تمامًا.
  \item هل $A(2, 1)$ و $B(5, 3)$ و $C(11, 7)$ على استقامة واحدة
        (\cref{met:g11:vect:alignment})؟
\end{enumerate}

\medskip
\noindent\textbf{الجزء الثاني --- خدمة حركة السفن.}
في اللحظة $t$ (بالساعات)، تكون السفينة $P$ في $(2t,\ 3 + t)$ وتكون السفينة $Q$
في $(10 - t,\ 2t - 1)$ (والمسافات بالأميال البحرية).
\begin{enumerate}[resume]
  \item أوجد المعادلتين الديكارتيتين \emph{للمسارين}،
        واحسب أين يتقاطعان.
  \item في أي لحظة تمر كل سفينة بنقطة \omterm{def:g10:numbers:interunion}{التقاطع}؟ وهل
        تتصادم السفينتان؟ صُغ التحذير العام في جملة واحدة:
        \omterm{def:g10:numbers:interunion}{تقاطع} المسارات في مواجهة التقاء الحركات.
  \item احسب مربع المسافة $D^2(t)$ بين السفينتين
        في اللحظة $t$، وردّه إلى مقدار تربيعي في $t$. وفي
        أي لحظة تكون السفينتان أقرب ما تكونان، وكم يبلغ اقترابهما
        (بتقنية رأس القطع في \cref{pb:g11:quad:1})؟
  \item تفرض الأنظمة فصلًا لا يقل عن \omterm{def:g10:reffunc:affine}{ميل} واحد.
        فهل يُنتهك؟ وإن كان كذلك، فحل $D^2(t) < 1$ وأعطِ
        مدة الانتهاك.
  \item فسّر لماذا يكون $D^2(t)$، من أجل \emph{أي} سفينتين على
        مسارين مستقيمين بسرعتين ثابتتين، \omterm{def:g11:quad:quadratic}{دالة تربيعية} في $t$ دائمًا ---
        فيكون ``أقرب نقطة اقتراب'' دائمًا
        حساب رأس.
  \item الاعتراض: يبدأ زورق دورية من المبدأ عند
        $t = 0$ بسرعة ثابتة $(a, b)$ ويجب أن يلتقي
        السفينة $P$ عند $t = 3$ بالضبط. احسب $(a, b)$ المطلوبة
        وسرعة زورق الدورية.
\end{enumerate}

\medskip
\noindent\textbf{الجزء الثالث --- \omterm{def:g10:coordgeom:system}{الإحداثيات} الذكية.}
اعمل في معلم \cref{exo:g11:vect:11}: فيصير متوازي الأضلاع
$ABCD$ هو $A(0,0)$ و $B(1,0)$ و $C(1,1)$ و $D(0,1)$؛ ولتكن $I$
\omterm{prop:g10:coordgeom:midpoint}{منتصف} $[AB]$ ولتكن $K$ \omterm{prop:g10:coordgeom:midpoint}{منتصف} $[CD]$.
\begin{enumerate}[resume]
  \item احسب \omterm{def:g10:algebra:equation}{معادلة} ديكارتية للمستقيم $(DI)$.
  \item \omterm{def:g11:deriv:rate}{قاطع} $(DI)$ مع القطر $(AC)$: واستعد
        واقعة \cref{exo:g11:vect:11} --- أي أن نقطة القطع تقع بالضبط
        على ثلث ارتفاع القطر.
  \item والآن المستقيم $(BK)$: أوجد معادلته، وقاطعه مع
        $(AC)$، واستنتج المبرهنة الكلاسيكية الكاملة:
        \emph{يقسم المستقيمان $(DI)$ و $(BK)$ القطر
        $(AC)$ ثلاثة أقسام متساوية}.
  \item لماذا كان مشروعًا برهان المبرهنة في هذا المعلم
        الخاص وحده؟ (ما الذي يحفظه اختيار \omterm{def:g10:coordgeom:system}{الإحداثيات}،
        وما الذي تتضمنه العبارة؟) أجب
        في جملتين أو ثلاث.
  \item وتثليث آخر للتمرن: لتكن $E$
        \omterm{prop:g10:coordgeom:midpoint}{منتصف} $[BC]$. أين يقطع المستقيم $(AE)$
        القطر \emph{الآخر} $(DB)$؟
\end{enumerate}

\medskip
\noindent\textbf{الجزء الرابع --- أحاجي المواضع.}
\begin{enumerate}[resume]
  \item هل المستقيمات الثلاثة $x + y = 4$ و \ $2x - y = 2$ و \
        $x - 2y = -2$ متلاقية؟ (\omterm{def:g11:deriv:rate}{قاطع} اثنين واختبر
        الثالث.)
  \item من أجل كل \omterm{def:g10:numbers:sets}{عدد حقيقي} $m$، ليكن $d_m$ المستقيم
        $mx - y + 2 - 3m = 0$. بيّن أن \emph{كل} $d_m$
        يمر \omterm{pb:g10:functions:1}{بنقطة ثابتة} واحدة. (اجمع الحدود التي فيها
        $m$.)
  \item في العائلة $(d_m)$: أي عضو يوازي
        $y = 2x + 7$؟ وأيها يمر بالمبدأ؟
  \item الخاتمة --- أثواب المستقيم الثلاثة: الديكارتي
        (اختبارات الانتماء بتعويض واحد)، والمختزل (\omterm{def:g10:reffunc:affine}{الميل}
        \omterm{def:g10:functions:graph}{والمنحنى} بنظرة)، والوسيطي (الحركة، والتوقيت، والاعتراض)؛
        \omterm{def:g11:vect:det}{والمحدد} عرّافًا للتوازي؛
        والمعلم المختار جيدًا مصنعًا للمبرهنات. جملة
        واحدة لكل بند.
\end{enumerate}
\end{problem}
